\section{Institutional Background and Related Work}\label{sec:background}

\paragraph*{The GENIUS Act} The Guiding and Establishing National Innovation for U.S. Stablecoins Act (S.\,1582) passed the Senate on 17~June 2025 and the House on 17~July 2025, and was signed on 18~July 2025 as Public Law 119-27~\cite{genius2025}. Section~4(a)(1)(A) requires a permitted issuer to hold identifiable reserves at least one-for-one, comprising US currency and Federal Reserve balances; demand deposits at insured institutions; ``Treasury bills, notes, or bonds (I)~with a remaining maturity of 93 days or less; or (II)~issued with a maturity of 93 days or less''; overnight repurchase and reverse repurchase agreements backed by Treasuries; government money market funds invested in these assets; other similarly liquid federal assets approved by the regulator; and tokenised forms of the above. The 93-day limit is therefore a ceiling on the maturity of eligible Treasuries, not a mandate to hold bills: an issuer can meet the rule entirely with deposits or repo.

Under \S20 the Act takes effect on the earlier of 18~months after enactment (18~January 2027) or 120~days after final implementing regulations. The Office of the Comptroller of the Currency issued its final rule on 21~October 2026 with an effective date of 18~January 2027~\cite{occ2026}. Throughout our post-enactment sample the reserve requirements were therefore known but not yet binding. ``Post'' in this paper denotes an \emph{anticipation} regime.

\paragraph*{Stablecoins and safe-asset prices} Ahmed and Aldasoro~\cite{ahmed2025} estimate the effect of stablecoin flows on Treasury yields with local projections~\cite{jorda2005}. The first version (May~2025) instruments flows with crypto-market shocks, reports that a two-standard-deviation inflow lowers 3-month yields by 2--2.5\,bp within ten days, and finds outflows raise yields two to three times as much as inflows lower them. The BIS figures quoted above (2.5--5\,bp per USD~3.5\,bn, up to three times larger in redemption episodes)~\cite{bis2025aer} predate the revision and draw on this version. The revised version (June~2026) uses a granular instrument across stablecoin--blockchain pairs, a sample from January~2021 to March~2026 and a larger set of controls. It reports $-0.71$\,bp on impact and about $-4$\,bp within ten days per USD~3.5\,bn. Its controls-only local projection, the specification closest to ours, gives less than 1\,bp within ten days and is not significant. The revision reports no asymmetry result and, in its footnote~25, does not use GENIUS passage as an event.

\paragraph*{Why maturity matters} In preferred-habitat models of the term structure~\cite{vayanos2021}, demand from investors anchored to a maturity segment moves yields in that segment most, with arbitrage transmitting part of the effect to neighbouring maturities. A rule confining stablecoin Treasury demand to $\leq$93~days thus predicts the pattern tested by~H2.
